% Section 7.6, from lectures/L07/README.md: the objectives, the questions to test yourself, and the
% next lecture.

\needspace{10\baselineskip}
\section{Review}
\label{c7:sec:review}

\subsection*{What you should be able to do now}
After this chapter, you should be able to:
\begin{itemize}
  \item Explain what separates a timer from a plain counter, and why the difference is one
    comparison.
  \item Build a timer as a gate network: an equality comparator from one XNOR per bit, and the
    clear that makes it repeat rather than wrap.
  \item Implement a timer in VHDL as a reusable module, converting a period in seconds into a
    \code{TICK_COUNT}, and say why \code{timeout} is a one-cycle pulse rather than a level.
  \item Compose a design out of pre-built modules rather than rewriting them, and say why a button
    controlling a timer still has to pass through a synchronizer first.
\end{itemize}

\needspace{6\baselineskip}
\subsection*{Questions to test yourself}
You should be able to explain:
\begin{enumerate}
  \item In what sense is a timer ``just a counter with a comparison''?
  \item Why must the internal counter be cleared back to \code{0} on reaching its target, rather
    than left to keep counting?
  \item Why is \code{timeout} a single-cycle pulse rather than a level that stays high?
  \item In the CircuitVerse version, changing the target count means rewiring the comparator. What
    replaces that rewiring in the VHDL version, and when is its value fixed?
  \item Why is it still necessary to run a button through a synchronizer before using its edge to
    enable a timer, even though the button has nothing to do with counting?
  \item \code{walking_led} writes almost no logic of its own. Name the three modules it composes
    and which chapter each came from.
\end{enumerate}

\needspace{6\baselineskip}
\subsection*{Looking ahead}
Chapter~\ref{c8:ch} is about state machines:
\begin{itemize}
  \item Finite state machines, the last building block in the book, designed by hand first: state
    diagram, state table, Karnaugh-derived logic, and a gate network in CircuitVerse, exactly the
    way this chapter's timer was drawn before it was written.
  \item Then in VHDL, where an enumerated type and a \code{case} statement replace the state table
    and the Karnaugh maps entirely.
  \item A timer is itself a state machine with two states and a single purpose.
    Chapter~\ref{c8:ch} generalizes that to circuits with many states and many transitions between
    them.
\end{itemize}
